%=============================================================
% Paper 2: Hopenhayn (1992), Econometrica
%=============================================================
\chapter{Hopenhayn (1992)}

\begin{paperbox}{Entry, Exit, and Firm Dynamics in Long Run Equilibrium}
\textbf{Author:} Hugo A.\ Hopenhayn\\[3pt]
\textbf{Journal:} \textit{Econometrica} 60(5): 1127--1150\\[3pt]
\textbf{Year:} 1992 \quad\textbf{Field:} Applied microeconomics; Industrial
organisation; Macroeconomics\\[3pt]
\textbf{Classification:} Theoretical
\end{paperbox}

%------------------------------------------------------------
\section{Research Question}

What determines the rates of firm entry, exit, and the stationary distribution
of firm size and profits in a competitive industry subject to idiosyncratic
productivity shocks, and how do changes in industry parameters (entry costs,
fixed costs, demand) affect these equilibrium outcomes?

%------------------------------------------------------------
\section{Motivation}

By the late 1980s, empirical work (Davis--Haltiwanger 1991; Dunne, Roberts,
Samuelson 1989) had documented two striking facts: (i) firm-specific
idiosyncratic uncertainty dominates firm-size dynamics; and (ii) industry-level
entry and exit rates are highly correlated, with large, persistent
cross-industry differences. Existing theoretical models either lacked firm
heterogeneity, lacked entry/exit in steady state, or were computationally
intractable dynamical systems. Hopenhayn builds a tractable equilibrium model
that generates ongoing entry and exit in the steady state, accommodates
idiosyncratic productivity shocks, admits closed-form comparative statics, and
can be connected to the growing empirical literature on firm heterogeneity.

%------------------------------------------------------------
\section{Main Contribution}

\begin{enumerate}
  \item \textbf{Stationary equilibrium concept:} Extends long-run competitive
    industry equilibrium to a setting with idiosyncratic shocks, firm
    heterogeneity, and simultaneous entry and exit in steady state. Proves
    existence and uniqueness under standard conditions (Theorems 2, 3, 4).

  \item \textbf{Comparative statics:} Derives clean predictions for how entry
    costs, fixed costs, and demand changes affect the exit threshold, firm
    turnover rate, and the stationary size distribution---including the
    counterintuitive result that higher entry costs \emph{protect incumbents}
    and \emph{reduce} turnover.

  \item \textbf{Lifecycle implications:} Shows that under monotone conditional
    dominance (mcd) ordering of the Markov transition, the distribution of
    firms' productivity shocks is stochastically increasing in cohort
    age---consistent with declining hazard rates, larger average size, and
    higher average profits for older cohorts (Propositions 3--5, 7--9).
\end{enumerate}

%------------------------------------------------------------
\section{Relationship to the Literature}

The paper extends and synthesises three strands:
\begin{itemize}
  \item \textbf{Jovanovic (1982):} passive learning model has no stationary
    equilibrium with entry/exit (the stopping time has positive probability at
    infinity). Hopenhayn's recurrence assumption (A.4) rules this out.
  \item \textbf{Lucas--Prescott (1971):} competitive equilibrium theory but
    with homogeneous firms and no entry/exit.
  \item \textbf{Ericson--Pakes (1989):} rich active-learning model but
    computationally intractable for analytical work.
\end{itemize}

The paper directly motivates Hopenhayn--Rogerson (1993) on firing costs,
Melitz (2003) on trade with heterogeneous firms, Asplund--Nocke (2006) on
imperfect competition, and Acemoglu et al.\ (2018) on innovation and
reallocation.

%------------------------------------------------------------
\section{Theoretical Framework and Economic Mechanism}

\subsection*{Industry Structure}

A continuum of firms produce a homogeneous good. Output market is given by
inverse demand $D(Q)$; input market (labour) by supply schedule $W(N)$.
Firms face a fixed cost $c_f$ per period and entry cost $c_e$ (sunk).
Individual firm output is $q=f(\varphi,n)$, where $\varphi \in [0,1]$ is an
idiosyncratic productivity shock evolving as a Markov process
$F(\varphi'|\varphi)$ satisfying first-order stochastic dominance in $\varphi$
(Assumption A.3(b)).

\subsection*{Firm's Bellman Equation}

The value of an incumbent firm with productivity $\varphi$ in a market with
distribution $\mu$ of active firms satisfies:
\begin{equation}
  v(\varphi,\mu) = \pi(\varphi,\mu)
    + \max\!\left\{0,\; \beta\int v(\varphi',\mu)\,F(d\varphi'|\varphi)\right\},
  \label{eq:hop_bellman}
\end{equation}
where $\pi(\varphi,\mu)$ is current-period profit (net of fixed cost $c_f$).
Because $v$ is strictly increasing in $\varphi$ (Proposition~3(i)), the
optimal exit rule is a reservation rule: exit if and only if
$\varphi < x^{*}$.

\subsection*{Equilibrium Conditions}

\textbf{Exit condition:}
\begin{equation}
  \int v(\varphi',\mu^*)\,F(d\varphi'|x^*) = 0.
  \label{eq:hop_exit}
\end{equation}
\textbf{Free entry condition:}
\begin{equation}
  \int v(\varphi,\mu^*)\,\nu(d\varphi) = c_e,
  \label{eq:hop_entry}
\end{equation}
where $\nu$ is the initial distribution from which entrants draw their
first productivity shock.

\textbf{Invariant measure:}
\begin{equation}
  \mu^* = P_x\,\mu^* + M^*\,\nu,
  \label{eq:hop_invariant}
\end{equation}
where $P_x$ is the transition operator conditional on survival ($\varphi \geq
x^*$) and $M^*$ is the equilibrium mass of entrants. The existence proof
constructs two curves $M_1(x)$ (exit-consistent entry mass) and $M_2(x)$
(entry-consistent entry mass) and establishes their intersection corresponds
to a stationary equilibrium with positive entry and exit (Theorem~3).

\subsection*{Lifecycle Mechanism}

As cohorts age, the worst-performing firms ($\varphi < x^*$) are continuously
removed by the exit rule. Under the mcd ordering of $F$, the distribution
$\mu_t$ of cohort-$t$ firms is stochastically increasing over time (Proposition~4),
delivering declining hazard rates, rising average size, and rising average
profits---all consistent with the empirical record.

%------------------------------------------------------------
\section{Data}

No empirical data. The paper is purely theoretical. Empirical facts motivating
the framework (approximately one-third of manufacturing jobs and over 40\% of
firms disappear over five-year periods) are drawn from Davis--Haltiwanger
(1991) and Dunne, Roberts, Samuelson (1989). A numerical example in Section~5
illustrates the comparative statics for the effect of fixed costs on the exit
threshold.

%------------------------------------------------------------
\section{Empirical Strategy}

Not applicable. The model's cross-sectional implications (declining hazard
rates with age, increasing average size with age, regression to the mean) are
compared informally to empirical evidence from Dunne, Roberts, and Samuelson.
No formal estimation or hypothesis testing is performed.

%------------------------------------------------------------
\section{Identification Strategy}

Not applicable in a causal sense. The paper's key analytical tool is the
algorithmic existence proof: showing that $x^*$ and $M^*$ are solutions to
the $M_1(x)$--$M_2(x)$ crossing, and that this crossing is unique under
Condition U.1 (price-taking in input markets) or U.2 (profit function
separable in $\varphi$ and prices). The proof is constructive and provides a
computational roadmap for numerical implementation.

%------------------------------------------------------------
\section{Main Results}

\subsection*{Existence and Uniqueness}

\begin{itemize}
  \item \textbf{Theorem 2:} A stationary competitive equilibrium always exists.
  \item \textbf{Theorem 3:} A stationary equilibrium with positive entry and
    exit exists if and only if $c_e < c^*$ for some threshold $c^*$.
  \item \textbf{Theorem 4:} Under Condition U.1 or U.2, the equilibrium is
    unique.
\end{itemize}

\subsection*{Lifecycle Properties}

Under the mcd ordering of $F$:
\begin{itemize}
  \item The age-cohort distributions $\{\mu_t\}$ are stochastically increasing
    in $t$ (Proposition~4).
  \item Hazard rates are declining in both age and current productivity
    (Proposition~5).
  \item Conditional on survival, older and larger firms have higher expected
    productivity, profits, and value.
\end{itemize}

\subsection*{Comparative Statics}

\begin{center}
\small
\begin{tabular}{p{3.5cm}p{3.5cm}p{3.5cm}p{2.5cm}}
\toprule
\textbf{Parameter change} & \textbf{Effect on $x^*$} &
\textbf{Effect on turnover rate} & \textbf{Condition} \\
\midrule
$\uparrow c_e$ & $\downarrow$ (less selection) & $\downarrow$ & Always \\
$\uparrow c_f$ & $\uparrow$ (more selection) & $\uparrow$ &
  Under U.2; can reverse under U.1 \\
Mean-preserving spread of $\nu$ & $\uparrow$ & $\uparrow$ &
  Under convexity of $v$ \\
$\uparrow$ Demand (elastic supply) & No effect & No effect &
  Under U.1 \\
\bottomrule
\end{tabular}
\end{center}

\subsection*{Profit and Value Bounds}

\textbf{Proposition 7:} Average industry profit satisfies
$\bar{\pi}(\mu^*) \geq \rho\, c_e$, where $\rho = M^*/\mu^*(S)$ is the
turnover rate. Turnover measures the minimum required average profit to
sustain the entry incentive.

\textbf{Proposition 8:} Average firm value $\bar{V}(\mu^*)=\bar{\pi}(\mu^*)/(1-\beta)
- \beta\rho c_e/(1-\beta)$, providing a formula for measuring implicit sunk
entry costs from observed value and profit data.

%------------------------------------------------------------
\section{Economic Magnitude and Interpretation}

The paper is analytical; magnitudes are structural rather than calibrated. The
key economic insight about magnitudes is the profit lower bound: an industry
with a 20\% annual turnover rate ($\rho=0.2$) and entry cost equal to one
year's revenue requires average annual profits of at least 20\% of entry cost
to sustain equilibrium entry. The numerical example in Section~5 illustrates
that the effect of higher fixed costs on the exit threshold can \emph{reverse
sign} depending on whether the technology features sufficiently high elasticity
of profits with respect to productivity---a subtle but important caveat on the
comparative statics.

%------------------------------------------------------------
\section{Mechanisms}

The central mechanism is \textbf{selection through exit}: the industry is in
a constant process of creative destruction. Each period, the worst-performing
firms ($\varphi < x^*$) exit, their resources are freed, and new entrants draw
from $\nu$. Three secondary mechanisms are worth noting:

\begin{itemize}
  \item \textbf{Option value of staying:} An incumbent with $\varphi$ slightly
    below $x^*$ may remain if the continuation value from potential future
    productivity improvement exceeds current negative profit. This is why
    $x^* < \varphi_{\max}$ and why even slightly unprofitable firms survive
    temporarily.
  \item \textbf{Selection amplification by age:} As cohorts age, selection
    removes the worst performers, leaving an increasingly productive and stable
    surviving set---consistent with empirically documented declining hazard rates.
  \item \textbf{Competition via entry and prices:} Higher entry rates drive
    down prices, raising the exit threshold and accelerating exit of less
    productive incumbents. This general equilibrium price mechanism is absent
    from partial equilibrium analyses.
\end{itemize}

%------------------------------------------------------------
\section{Robustness Checks}

As a theoretical paper, robustness is assessed through:
\begin{itemize}
  \item Proof of existence under broad conditions (Assumptions A.1--A.5 only).
  \item Sufficient conditions for uniqueness (U.1 or U.2, both widely
    satisfied by standard functional forms including CES and Cobb--Douglas).
  \item A numerical example demonstrating that the fixed-cost comparative
    static can reverse under U.1, providing an important caveat.
  \item Extension to multidimensional state spaces (CES production with
    multiple productivity parameters) shown to satisfy U.2.
\end{itemize}

%------------------------------------------------------------
\section{Limitations}

\begin{enumerate}
  \item \textbf{No dynamics:} Focuses entirely on stationary equilibria.
    Transitional dynamics---how the industry responds to a parameter change
    before reaching the new stationary state---are not analysed.
  \item \textbf{Competitive markets only:} Extension to imperfect competition
    requires Asplund--Nocke (2006); price-competition effects are absent.
  \item \textbf{No aggregate shocks:} All aggregate variables (price, output,
    employment) are deterministic in the stationary state. There are no
    business cycles.
  \item \textbf{Passive productivity process:} $\varphi$ evolves exogenously.
    Firms cannot invest to change their position in the productivity
    distribution (handled by Ericson--Pakes and Acemoglu et al.).
  \item \textbf{No worker side:} Labour is supplied elastically at price
    $W(N)$. Labour market frictions, search, or matching are outside the model.
\end{enumerate}

%------------------------------------------------------------
\section{Policy Implications}

\begin{itemize}
  \item \textbf{Entry barriers reduce turnover and protect incumbents:} Higher
    $c_e$ reduces $x^*$, enabling less productive firms to survive. Entry
    barriers impose a double welfare cost: fewer entrants and a less efficient
    incumbent stock.
  \item \textbf{Fixed costs and selection:} Higher fixed production costs
    raise $x^*$, increasing average productivity by forcing exit of less
    productive firms. This provides a theoretical basis for why fixed-cost
    regulation can have ambiguous welfare effects.
  \item \textbf{Firing cost policies:} The paper (final remarks) notes that
    a version of the model is used in Hopenhayn--Rogerson (1993) to study
    firing costs. By raising effective fixed costs and reducing selection,
    firing costs reduce long-run productivity.
\end{itemize}

%------------------------------------------------------------
\section{Overall Assessment}

Hopenhayn (1992) is one of the most important theory papers in industrial
organisation and macroeconomics of the last four decades. Its significance
lies not in the complexity of the model---which is deliberately
parsimonious---but in the analytical tractability it achieves while
incorporating heterogeneity, entry, and exit simultaneously. The existence--
uniqueness results, comparative statics, and lifecycle propositions have all
become standard ingredients in quantitative macro models. Direct descendants
include Melitz (2003), Hopenhayn--Rogerson (1993), and Acemoglu et al.\
(2018). The paper's limitations are largely features: by keeping the model
simple, Hopenhayn made it usable.

\textbf{Quality: Essential. A must-read for anyone working on firm dynamics,
industrial policy, or models with heterogeneous agents.}

%------------------------------------------------------------
\section{Relevance to My Research}

Directly relevant as the theoretical foundation for any RDC project on firm
entry, exit, and dynamics. The comparative statics on entry costs and fixed
costs provide testable predictions for Canadian data: industries with lower
entry barriers should show higher turnover rates and higher productivity
distributions. The profit bound $\bar{\pi} \geq \rho\, c_e$ provides a
formula evaluable with Canadian firm-level data (T2 corporate tax records).

%------------------------------------------------------------
\section{Potential Research Extensions}

\begin{itemize}
  \item Estimate the model structurally using Canadian firm-level data to
    recover $c_e$, $c_f$, and the Markov process $F(\varphi'|\varphi)$.
  \item Exploit provincial variation in effective entry costs (regulatory
    burden indices) to test the comparative statics predictions on firm
    turnover and productivity distributions.
  \item Embed in a GE model with labour market search to study the interaction
    between firm exit rates and unemployment duration.
\end{itemize}

%------------------------------------------------------------
\section{Research Gaps}

\begin{itemize}
  \item \textbf{Transitional dynamics:} Policy reform evaluations require
    knowing the transition path, which the stationary analysis cannot provide.
  \item \textbf{Worker heterogeneity:} Workers are homogeneous; a natural
    extension would allow differential matching with firm types, generating
    wage variation driven by firm productivity heterogeneity.
  \item \textbf{Financial constraints:} Introducing credit constraints could
    generate amplification---firms below the exit threshold that would survive
    absent the constraint, but cannot access external finance.
\end{itemize}

%------------------------------------------------------------
\section{Methodological Lessons}

\begin{enumerate}
  \item The stationary equilibrium concept---requiring that the invariant
    measure $\mu^*$ equals the measure generated by entry/exit rules and
    entry distribution---is the right way to think about long-run industry
    equilibria with heterogeneous firms.
  \item The algorithmic existence proof (constructing $M_1(x)$ and $M_2(x)$
    curves) provides a computational roadmap for numerical implementation---
    a rare feature of analytical theory papers.
  \item The mcd (monotone conditional dominance) ordering is stronger than
    FOSD but weaker than likelihood ratio ordering, and survives the truncation
    from the exit threshold---a useful technical tool for dynamic selection
    models.
\end{enumerate}

%------------------------------------------------------------
\section*{Five-Sentence Summary}

Hopenhayn (1992) develops a tractable dynamic stochastic model of a
competitive industry in which firms face idiosyncratic productivity shocks and
choose optimally when to exit, with free entry driving expected profits to the
entry cost. The concept of stationary competitive equilibrium---a steady-state
where entry and exit rates are equal and the distribution of active firms is
time-invariant---is introduced and shown to exist and be unique under broad
conditions. Comparative statics demonstrate that higher entry costs reduce
turnover by protecting less productive incumbents, while higher fixed costs
(under separability) raise the exit threshold and improve average firm
productivity through selection. Under monotone conditional dominance of the
productivity process, the model implies that older cohorts of firms are
stochastically larger, have lower hazard rates, and earn higher average
profits---all consistent with the empirical evidence from Dunne, Roberts, and
Samuelson. The paper provides the canonical theoretical framework for all
subsequent quantitative work on firm dynamics, entry, exit, and the long-run
effects of industrial policy on productivity and resource allocation.

\vspace{6pt}
\begin{quickref}
\begin{tabularx}{\linewidth}{@{}lX@{}}
\textbf{Key contribution} &
  Tractable, analytically characterised stationary equilibrium with
  simultaneous entry and exit driven by idiosyncratic productivity; clean
  comparative statics on entry/fixed costs and demand. \\[4pt]
\textbf{Key weakness} &
  No aggregate shocks (no business cycle); no transitional dynamics;
  perfectly competitive markets only; no investment in productivity. \\[4pt]
\textbf{Most interesting gap} &
  Transitional dynamics after policy changes---a reform that reduces entry
  costs generates extended transitional dynamics that the stationary analysis
  cannot characterise. \\[4pt]
\textbf{Publishable extension} &
  Structural estimation of the Hopenhayn model using Canadian T2 tax data,
  exploiting variation in effective provincial entry costs (regulatory burden
  indices) to identify the entry cost parameter and test comparative statics
  predictions on firm turnover and productivity distributions.
\end{tabularx}
\end{quickref}
